\chapter{MỞ ĐẦU VÀ TỔNG QUAN ĐỀ TÀI}

\section{Lý do chọn đề tài}

Hóa đơn chứa tên cửa hàng, ngày mua, danh sách mặt hàng và tổng tiền với bố cục thay đổi giữa các mẫu. Một số tiền có thể là đơn giá, thuế hoặc tổng thanh toán. Vì vậy, đọc được chữ chưa đủ; hệ thống còn phải xác định ý nghĩa của từng đoạn chữ.

Nhóm khảo sát ba hướng: chỉ dùng văn bản; dùng văn bản cùng vị trí; kết hợp văn bản, vị trí và đặc trưng ảnh. Thực nghiệm trên cùng dữ liệu giúp so sánh hiệu quả của các cấu hình này.

\section{Mục tiêu và phạm vi}

Mục tiêu là khảo sát, triển khai và so sánh các mô hình trong điều kiện thực nghiệm thống nhất bằng Precision, Recall, F1. Báo cáo trình bày RoBERTa, BROS và PICK trên SROIE, CORD-v2 và WildReceipt, tương ứng chín checkpoint cho chín cặp mô hình--bộ dữ liệu.

Thực nghiệm chính dùng text và box từ chú thích gốc. Chương trình chuyển đổi đầu vào nhưng giữ nguyên dữ liệu nguồn.

\section{Nội dung đã hoàn thành}

Nhóm đã kiểm tra dữ liệu, chạy chín cấu hình và lưu checkpoint tốt nhất theo validation. Dự đoán theo ảnh và theo đoạn chữ được xuất lại trên Kaggle để xem lỗi, không huấn luyện lại. Báo cáo gồm năm chương, từ tổng quan và cơ sở lý thuyết đến triển khai, thực nghiệm và kết luận; số liệu được cập nhật đến ngày 24/10/2026.
